\documentclass[10pt,a4paper]{amsart}
\usepackage[margin=19mm]{geometry}
\usepackage{amsmath,amssymb,mathtools}
\usepackage[scr=rsfs,cal=euler]{mathalfa}
\usepackage{xfrac}
\usepackage[unicode,hidelinks,bookmarks=false]{hyperref}
\newcommand{\ie}{i.e.\ }
\newcommand{\cf}{cf.\ }
\newcommand{\resp}{respectively\ }
\newcommand{\Subsection}{\subsection}
\providecommand{\nobreakdash}{\nobreak\hbox{-}}
\setlength{\emergencystretch}{2em}
\multlinegap0pt
\usepackage{mathtools}
\usepackage{amssymb}
\usepackage[shortlabels]{enumitem}
\newtheorem{theorem}{Theorem}
\title{results on dynamics of bungee set of composite entire functions in the eremenko-lyubich class}
\author{Dinesh Kumar, Soumyajeet Das}
\date{Source: arXiv:2405.11217v4 (CC BY 4.0)}
\begin{document}
\maketitle

\noindent\emph{Source and licence.} results on dynamics of bungee set of composite entire functions in the eremenko-lyubich class. Version arXiv:2405.11217v4 (CC BY 4.0), distributed by arXiv under the Creative Commons Attribution 4.0 International licence, http://creativecommons.org/licenses/by/4.0/, as recorded in the arXiv licence metadata for this exact version and shown on the abstract page https://arxiv.org/abs/2405.11217v4. Full licence notice: \texttt{licenses/arxiv-2405.11217-CC-BY-4.0.txt}.\par\smallskip

\noindent\textit{Editorial scope.} Counted unit: the article's theorem theorem (source0.tex lines 292--294). The article's complete original proof of it is delivered with it (source0.tex lines 296--312, one \texttt{proof} environment of 17 lines). No mathematics is written, repaired or re-derived: the statement and the proof are the article's own lines, in the article's own notation, and no retained line is substituted or re-flowed.  Retained uncounted as the source-local context the proof needs: lines 215--217.  Nothing was omitted from any retained span, so the package carries no omission record.  No retained line contains a citation, so the wrapper carries no bibliography and no bibliography entry.  No retained line contains a figure environment, an image-inclusion command or drawing code, so no image is delivered and the package declares no asset dependency.  Editorial formatting only, in addition: an AMS \texttt{amsart} preamble that restates the article's own declarations and macros used by the retained lines, namely the environments \texttt{theorem} on separate counters, together with the standard packages those lines need.  The retained environments print in this document as Theorem 1 in order of appearance, whereas the article prints the same items with the numbering of its own full document.  The complete native source of the article as distributed in the arXiv e-print for this exact version is bundled unchanged as \texttt{sources/159352/source0.tex}.
\begin{theorem} \label{Complete_invar_esp}
If $f$ and $g$ are commuting functions satisfying \textbf{Property A}, then $I(f)$ is completely invariant under $g$ and vice versa. 
\end{theorem}

\begin{theorem}
Suppose $f$ and $g$ are commuting transcendental entire functions  satisfying \textbf{Property A}. Then, $BU(f \circ g)\subset BU(f)\cap BU(g).$
\end{theorem}

\begin{proof}
%We  establish that $BU(f \circ g)\subset BU(f)\cap BU(g).$

%Let $f$ and $g$ be two permutable entire functions such that both $f$ and $g$ do not have any finite asymptotic values. Then $BU(f\circ g)\subset BU(f)\cap BU(g)\subset BU(f)\cup BU(g)$.

 Suppose $z_0\in BU(f\circ g).$ Then,  there exists two subsequences $\{n_k\},\{m_k\}$ and a constant $M>0$ such that $(f\circ g)^{n_k}(z_0)\to\infty$ as $k\to\infty$ and $|(f\circ g)^{m_k}(z_0)|<M \mbox{ for all } k\in \mathbb{N}$. It is enough to show that $z_0\in BU(g)$ (as the proof of $z_0\in BU(f)$ follows on similar lines). In order to show  this, we prove that $z_0$ cannot be in $I(g)$ as well as $K(g).$ We divide the proof into several cases:\\
Case(i): Suppose that $g^{n_k}(z_0)\to\infty$ and $g^{m_k}(z_0)\to\infty$ as $k\to\infty$.\\
We can assume that $g^n(z_0)\to\infty$ as $n\to\infty$
(for, if there exists a subsequence $\{l_k\}$ such 
that $g^{l_k}(z_0)$ stays bounded as $k\to\infty$ then 
this will clearly imply that $z_0\in BU(g)$). Also, using Theorem \ref{Complete_invar_esp}, $f^k(I(g)\subset I(g)$ for every  $k\in\mathbb{N}$. In particular, $f^{m_k}(z_0)\in I(g)$ which implies that $(g\circ f)^{m_k}(z_0)\to\infty $ as $k\to\infty$ which leads to a  contradiction. \\
Case(ii): Suppose that $g^{n_k}(z_0)\to b$ and $g^{m_k}(z_0)\to\infty$ as $k\to\infty$. This again shows that $z_0\in BU(g)$ and the result follows.\\
Case(3): Suppose that $g^{n_k}(z_0)\to \alpha $ and $g^{m_k}(z_0)\to \beta$ as $k\to\infty.$   It follows that $z_0\in K(g)$ (for,  if there exists a subsequence $\{l_k\}$ for 
which $g^{l_k}(z_0)\to\infty$ as $k\to\infty$ then we are back to Case(2) 
 and hence $z_0\in BU(g)$). Using Theorem \ref{Complete_invar_esp},  $f(K(g)\subset K(g)$. In particular, $f^{n_k}(z_0)\in K(g)$, i.e., there exists $L>0$ such that $|g^n(f^{n_k}(z_0))|<L \mbox{ for all } n\in\mathbb{N}$. This implies that $|(f\circ g)^{n_k}(z_0)|<L \mbox{ for all } k\in\mathbb{N}$ which is again a contradiction. Thus, from the above cases, we conclude that $z_0\in BU(g).$  The proof of $z_0\in BU(f)$ follows on similar lines. 
Hence, we obtain that $BU(f\circ g)\subset BU(f)\cap BU(g)\subset BU(f)\cup BU(g)$.
\end{proof}

\end{document}
